\section{Derivación de funciones complejas (de variable compleja)}

\subsection{Definición de derivada compleja}

Si $A \subset \mathbb{C}$ y $f : A \to \mathbb{C}$, podemos escribir $f$ como:
\[
    f(z) = f(x+iy) = u(x,y) + iv(x,y),
\]
donde $u(x,y) = \Re(f(x+iy))$, y $v(x,y) = \Im(f(x+iy))$ son funciones reales de dos variables.
Así podemos considerar $f$ como una función de $\mathbb{R}^2$ en $\mathbb{R}^2$:
\[
    f: A \subset \mathbb{R}^2 \to \mathbb{R}^2, \quad f(x,y) = (u(x,y), v(x,y)).
\]

\begin{definición}[Derivada compleja]
    Sea $\Omega \subset \mathbb{C}$ un abierto, $f : \Omega \to \mathbb{C}$ y $z_0 \in \Omega$.\\
    Diremos que $f$ es \textbf{derivable} en $z_0$ si existe el siguiente límite en $\mathbb{C}$:
    \[
        f'(z_0) = \lim_{z \to z_0} \frac{f(z) - f(z_0)}{z - z_0} = \lim_{h \to 0} \frac{f(z_0 + h) - f(z_0)}{h} = w \in \mathbb{C}.
    \]
    Es decir, esta derivada vale $w$ si para cada $\varepsilon > 0$ existe $\delta > 0$ tal que si $z \in \Omega$ y $0 < |z-z_0| < \delta$, entonces
    \[
        \left| \frac{f(z) - f(z_0)}{z - z_0} - w \right| < \varepsilon.
    \]
    A este limite, cuando existe, lo llamamos la \textbf{derivada} de $f$ en $z_0$ y lo denotamos por:
    \[
        f'(z_0) = \frac{df}{dz}(z_0) = \frac{d}{dz} f(z_0).
    \]
\end{definición}

Sea $f = u + iv$ una función compleja definida en un entorno de $z_0 = x_0 + i y_0$. 
Supongamos que $f$ es derivable en $z_0$ y que
\[
f'(z_0) = \alpha + i \beta, \qquad \alpha, \beta \in \mathbb{R}.
\]
Entonces se cumple la caracterización siguiente.

Por definición de derivada compleja,
\[
\lim_{h \to 0} \frac{f(z_0 + h) - f(z_0)}{h} = f'(z_0) = \alpha + i \beta.
\]

Esto equivale a exigir que
\[
\left| \frac{f(z_0 + h) - f(z_0)}{h} - (\alpha + i \beta) \right| \xrightarrow{h \to 0} 0,
\]
lo cual, separando partes real e imaginaria, se traduce en
\[
\begin{cases}
\Re\!\left( \dfrac{f(z_0 + h) - f(z_0)}{h} \right) \xrightarrow{h \to 0} \alpha, \\[1.2ex]
\Im\!\left( \dfrac{f(z_0 + h) - f(z_0)}{h} \right) \xrightarrow{h \to 0} \beta.
\end{cases}
\]

Ahora escribimos $h = h_1 + i h_2$, con $(h_1,h_2)\in\mathbb{R}^2$. Entonces:
\[
\begin{cases}
\dfrac{1}{|h|}\,\big|u(x_0+h_1,y_0+h_2) - u(x_0,y_0) - \Re[(h_1+i h_2)(\alpha+i\beta)]\big| \xrightarrow{|h|\to 0} 0, \\[2ex]
\dfrac{1}{|h|}\,\big|v(x_0+h_1,y_0+h_2) - v(x_0,y_0) - \Im[(h_1+i h_2)(\alpha+i\beta)]\big| \xrightarrow{|h|\to 0} 0.
\end{cases}
\]

Desarrollando el producto:
\[
(h_1+i h_2)(\alpha+i\beta) = (h_1\alpha - h_2\beta) + i(h_1\beta + h_2\alpha).
\]

Por tanto, las condiciones se reescriben como
\[
\begin{cases}
\dfrac{1}{\|(h_1,h_2)\|}\,\big|u(x_0+h_1,y_0+h_2) - u(x_0,y_0) - (h_1\alpha - h_2\beta)\big| \xrightarrow{(h_1,h_2)\to (0,0)} 0, \\[2ex]
\dfrac{1}{\|(h_1,h_2)\|}\,\big|v(x_0+h_1,y_0+h_2) - v(x_0,y_0) - (h_1\beta + h_2\alpha)\big| \xrightarrow{(h_1,h_2)\to (0,0)} 0.
\end{cases}
\]

Esto significa exactamente que:
\[
Du(x_0,y_0) = (\alpha,-\beta), \qquad Dv(x_0,y_0) = (\beta,\alpha),
\]
es decir, que $u$ y $v$ son diferenciables en $(x_0,y_0)$ y sus derivadas parciales satisfacen
\[
\frac{\partial u}{\partial x}(x_0,y_0) = \alpha, 
\qquad 
\frac{\partial u}{\partial y}(x_0,y_0) = -\beta,
\]
\[
\frac{\partial v}{\partial x}(x_0,y_0) = \beta, 
\qquad 
\frac{\partial v}{\partial y}(x_0,y_0) = \alpha.
\]

Por tanto,
\[
\frac{\partial u}{\partial x}(x_0,y_0) = \frac{\partial v}{\partial y}(x_0,y_0), 
\qquad 
\frac{\partial u}{\partial y}(x_0,y_0) = -\frac{\partial v}{\partial x}(x_0,y_0).
\]

Estas son precisamente las \textbf{ecuaciones de Cauchy–Riemann}.
\begin{teorema}
    Sea $\Omega \subset \mathbb{C}$ un abierto, $f : \Omega \to \mathbb{C}$ y $z_0 = x_0 + i y_0 \in \Omega$. Entonces, $f = u+iv$ es derivable en $z_0$ si y sólo si $u, v$ son diferenciables en $(x_0, y_0)$ y se cumplen las \textbf{ecuaciones de Cauchy-Riemann}:
    \[
        \frac{\partial u}{\partial x}(x_0, y_0) = \frac{\partial v}{\partial y}(x_0, y_0), \quad \frac{\partial u}{\partial y}(x_0, y_0) = -\frac{\partial v}{\partial x}(x_0, y_0).
    \]
    Y en tal caso, la derivada de $f$ en $z_0$ viene dada por:
    \[
        f'(z_0) = u_x(x_0, y_0) + i v_x(x_0, y_0) = v_y(x_0, y_0) - i u_y(x_0, y_0).
    \]
\end{teorema}

\begin{observación}
    Por el teorema anterior, $f : \Omega \to \mathbb{C}$ es derivable en $z_0 \in \Omega$ si y sólo si
    \[
        f : \Omega \subset \mathbb{R}^2 \to \mathbb{R}^2, \quad f(x,y) = (u(x,y), v(x,y))
    \]
    es diferenciable en $(x_0, y_0)$ y su matriz jacobiana es de la forma:
    \[
        J_f(x_0, y_0) = \begin{pmatrix}
            \alpha & -\beta \\
            \beta & \alpha 
        \end{pmatrix}, \quad \alpha, \beta \in \mathbb{R}.
    \]
\end{observación}

\ejemplo{
    ¿Dónde es derivable $f(z) = |z| \ \forall z \in \mathbb{C}$?\\
    \[
        u(z) = u(x,y) = \Re(f(z)) = \Re(|z|) = \sqrt{x^2 + y^2}, \quad v(z) = v(x,y) = \Im(f(z)) = \Im(|z|) = 0.
    \]
    \[
        \begin{cases}
            u(x,y) = \sqrt{x^2 + y^2} \\
            v(x,y) = 0
        \end{cases} \implies
        \begin{cases}
            u_x(x,y) = \frac{2x}{2\sqrt{x^2 + y^2}} = \frac{x}{\sqrt{x^2 + y^2}} \\
            u_y(x,y) = \frac{2y}{2\sqrt{x^2 + y^2}} = \frac{y}{\sqrt{x^2 + y^2}} \\
            v_x \equiv v_y \text{ en todo } \mathbb{R}^2
        \end{cases} \quad \forall (x,y) \in \mathbb{R}^2 \setminus \{(0,0)\}
    \]
    Si $z = x + iy \neq 0$, las ecuaciones de Cauchy-Riemann no se cumplen:
    \[
        \begin{cases}
            u_x(x,y) = v_y(x,y) \\
            u_y(x,y) = -v_x(x,y)
        \end{cases} \iff
        \begin{cases}
            \frac{x}{\sqrt{x^2 + y^2}} = 0 \\
            \frac{y}{\sqrt{x^2 + y^2}} = 0
        \end{cases} \iff (x,y) = (0,0) \text{ (imposible)}
    \]
    Por lo tanto, $f$ no es derivable en ningún punto $z \neq 0$.\\ 
    En $z=0$, tampoco lo es pues
    \[
        \lim_{z \to 0} \frac{f(z) - f(0)}{z - 0} = \lim_{z \to 0} \frac{|z| - 0}{z}
    \]
    y este límite no existe 
    \[
        \frac{1}{n} \xrightarrow{n \to \infty} 0, \quad \frac{|\frac{1}{n}|}{\frac{1}{n}} = 1 \xrightarrow{n \to \infty} 1.
    \]
    y
    \[
        -\frac{1}{n} \xrightarrow{n \to \infty} 0, \quad \frac{|\frac{-1}{n}|}{\frac{-1}{n}} = -1 \xrightarrow{n \to \infty} -1.
    \]
    Por lo tanto, $f$ no es derivable en ningún punto de $\mathbb{C}$.  
}

Se puede probar de forma análoga al caso real que

\begin{itemize}
    \item Si $f$ es derivable en $z_0$, entonces $f$ es continua en $z_0$.
    \item Si $f$ y $g$ son derivables en $z_0$, entonces $f+g$, $\alpha f$ ($\alpha \in \mathbb{C}$) y $f \cdot g$ son derivables en $z_0$, y se cumplen las reglas usuales de derivación, es decir,
    \[
        (f+g)'(z_0) = f'(z_0) + g'(z_0), \quad (\alpha f)'(z_0) = \alpha f'(z_0),
    \]
    \[
        (f \cdot g)'(z_0) = f'(z_0) g(z_0) + f(z_0) g'(z_0)
    \]
    \item Si $f$ y $g$ son derivables en $z_0$ y $g(z_0) \neq 0$, entonces $f/g$ es derivable en $z_0$ y
    \[
        \left(\frac{f}{g}\right)'(z_0) = \frac{f'(z_0) g(z_0) - f(z_0) g'(z_0)}{(g(z_0))^2}.
    \]
    \item Si $f$ es derivable en $z_0$ y $g$ es derivable en $f(z_0)$, entonces $g \circ f$ es derivable en $z_0$ y
    \[
        (g \circ f)'(z_0) = g'(f(z_0)) \cdot f'(z_0).
    \]
\end{itemize}

\begin{definición} [Funciones holomorfas y enteras]
    Sea $\Omega \subset \mathbb{C}$ un abierto y $f : \Omega \to \mathbb{C}$. Diremos que $f$ es \textbf{holomorfa} en $\Omega$ (y escribiremos $f \in \mathcal{H}(\Omega)$) si es derivable en todo punto de $\Omega$. Si $f$ es holomorfa en todo $\mathbb{C}$, diremos que $f$ es una función \textbf{entera}.
\end{definición}

\textbf{Notas:}

\begin{itemize}
    \item La función $f(z) = z |z|^2$ es derivable sólo en $z=0$.
    \item La función 
    \[
        f(z) = \begin{cases}
            \frac{(1+i)x^3 - (1-i)y^3}{x^2 + y^2}, \quad \forall z = (x,y) \in \mathbb{C} \setminus \{0\} \\
            0, \quad z=0
        \end{cases}
    \]
    cumple las ecuaciones de Cauchy-Riemann en $z=0$ pero no es derivable en $z=0$.
\end{itemize}

\begin{proposición}
    Sea $\Omega \subset \mathbb{C}$ un abierto conexo y $f : \Omega \to \mathbb{C}$ holomorfa. Entonces se tiene
    \begin{enumerate}
        \item Si $f'(z) = 0$ para todo $z \in \Omega$, entonces $f$ es constante en $\Omega$.
        \item Si $|f(z)|$ es constante en $\Omega$, entonces $f$ es constante en $\Omega$.
        \item Si $f(z) \in \mathbb{R}$ para todo $z \in \Omega$, entonces $f$ es constante en $\Omega$.
        \item Si $f(z) \in i \mathbb{R} = \{iw : w \in \mathbb{R}\}$ para todo $z \in \Omega$, entonces $f$ es constante en $\Omega$.
    \end{enumerate} \label{prop:constantes_holomorfas}
\end{proposición}

\begin{proof}
    \leavevmode
    \begin{enumerate}
        \item Sea $f = u + iv$, derivable en $\Omega$. Entonces $u, v$ son diferenciables en $\Omega$ y 
        \[
            f'(z) = 0 \implies u_x(z) = u_y(z) = v_x(z) = v_y(z) = 0 \quad \forall z = (x,y) \in \Omega.
        \]
        \[
            \implies u, v \text{ son constantes en } \Omega \implies f \text{ es constante en } \Omega.
        \]
        \item Sea $c \geq 0$ tal que $|f(z)| = c$ para todo $z \in \Omega$. Entonces
        \begin{itemize}
            \item Si $c=0$, entonces $f(z) = 0$ para todo $z \in \Omega$ y $f$ es constante en $\Omega$.
            \item Si $c \neq 0$ y llamamos 
            \[
                \begin{cases}
                    u(z) = (\Re(f(z))) \\
                    v(z) = (\Im(f(z)))
                \end{cases}
            \]
            tenemos que 
            \[
                c = |f(z)| = |u(z) + iv(z)| \ \forall z \in \Omega \implies u^2(z) + v^2(z) = c^2 \ \forall z \in \Omega.
            \]
            \[
                \implies \begin{cases}
                    2u(z) u_x(z) + 2v(z) v_x(z) = 0 \\
                    2u(z) u_y(z) + 2v(z) v_y(z) = 0
                \end{cases} \quad \forall z \in \Omega. \implies
                \begin{cases}
                    u(z) u_x(z) + v(z) v_x(z) = 0 \\
                    u(z) u_y(z) + v(z) v_y(z) = 0
                \end{cases} \quad \forall z \in \Omega.
            \]
            luego $f$ es holomorfa en $\Omega$ y se cumplen las ecuaciones de Cauchy-Riemann, por lo que
            \[
                \begin{cases}
                    u(z) u_x(z) - v(z) u_y(z) = 0 \\
                    u(z) u_y(z) + v(z) u_x(z) = 0
                \end{cases} \implies 
                \begin{cases}
                    u(z)^2 u_x(z) - u(z) v(z) u_y(z) = 0 \\
                    u(z) v(z) u_y(z) + v(z)^2 u_x(z) = 0 
                \end{cases}
            \]
            \[
                \implies u_x(z) (u^2(z) + v^2(z)) = 0 \implies u_x(z) c^2 = 0 \underbrace{\implies}_{c \neq 0} u_x(z) = 0 \quad \forall z \in \Omega.
            \]
            Análogamente, $v_x(z) = 0$ para todo $z \in \Omega$. Por lo tanto, $f'(z) = 0$ para todo $z \in \Omega$ y $f$ es constante en $\Omega$ por (1).
        \end{itemize}
        \item Si $f(z) \in \mathbb{R}$ para todo $z \in \Omega$, entonces 
        \[
            \begin{cases}
                u(z) = (\Re(f(z))) = f(z) \\
                v(z) = (\Im(f(z))) = 0
            \end{cases} \quad \forall z \in \Omega.
        \]
        \[
            \implies v_x(z) = v_y(z) = 0 \quad \forall z \in \Omega \implies f'(z) = 0 \quad \forall z \in \Omega \implies f \text{ es constante en } \Omega \text{ por (1).}
        \]
    \end{enumerate}

\end{proof}

Si $f$ es derivable en $z_0 \in \Omega$ con $\Omega \subset \mathbb{C}$ abierto, entonces tenemos que

\[
    \lim_{h \to 0} \frac{f(z_0 + h) - f(z_0)}{h} = f'(z_0) \in \mathbb{C}.
\]

Luego $f(z_0 + h) = f(z_0) + f'(z_0) h + h \varphi(h)$, donde $\varphi(h) \to 0$ cuando $h \to 0$. Por tanto, basta tomar 

\[
    \varphi(h) = \frac{f(z_0 + h) - f(z_0)}{h} - f'(z_0)
\]

Más aún, $f$ es derivable en $z_0$ y $f'(z_0) = \alpha + i \beta$ si y sólo $\exists \varphi : D(0,r_0) \to \mathbb{C}$ con $\lim_{h \to 0} \varphi(h) = 0$ tal que

\[
    f(z_0 + h) = f(z_0) + (\alpha + i \beta) h + h \varphi(h), \quad \forall h \in D(0,r_0).
\]

